Intracranial aneurysms are a frequent cerebrovascular abnormality; rupture may lead to subarachnoid hemorrhage with high morbidity and mortality. CTA depicts parent vessels and aneurysm morphology well before intervention, yet reviewing full volumetric series remains time-consuming, and inter-observer disagreement persists under complex vascular geometry and artifacts. Rupture risk is governed by heterogeneous clinical and morphologic factors; classical scales or shallow models struggle to capture nonlinear interactions between numeric and categorical signals. Building a computer-aided pipeline that robustly segments aneurysms on CTA and stratifies rupture risk is therefore clinically motivated.\par

This thesis implements two reproducible tracks aligned with the course project. For segmentation, a 3D Attention U-Net is adopted: an encoder--decoder backbone retains multi-scale skip connections, while attention gates re-weight skip features at each upsampling stage so that the network emphasizes plausible aneurysm regions in predominantly background voxels. Training uses case-wise streaming and patch-wise sampling with spacing resampling and intensity normalization; the objective combines binary cross-entropy with Dice loss to balance pixel-wise calibration and overlap. At inference time, sequential patch prediction is stitched back to a full volume and exported with restored image geometry for downstream use.\par

For rupture risk, tabular clinical and derived features are encoded in a dedicated RiskCrossAttentionModel: numeric inputs are scaled and projected per feature into tokens, categorical fields use embeddings, and stacked cross-attention blocks mutually refine the two streams before a fused MLP head outputs a rupture probability. The tabular pipeline covers imputation, rare-category pooling, mutual-information feature selection, and class-frequency--weighted loss; a validation-set threshold search optimizes F1 before fixing the operating point for testing.\par

Configuration is centralized in YAML; \texttt{train\_split.py} and \texttt{train\_risk.py} produce checkpoints and JSON metric summaries respectively. A lightweight FastAPI UI demonstrates volume upload, patch-based segmentation, and tabular risk inference. Quantitative claims in this abstract are intentionally conservative: final numbers should be taken from institution-approved datasets and the project's own evaluation logs after replication.\par

~\\
\hspace*{2em}\textbf{Key words}: intracranial aneurysm; CTA; 3D segmentation; Attention U-Net; tabular learning; cross-attention; rupture risk prediction\\
